%%
%% This is file `sample-manuscript.tex',
%% generated with the docstrip utility.
%%
%% The original source files were:
%%
%% samples.dtx  (with options: `manuscript')
%% 
%% IMPORTANT NOTICE:
%% 
%% For the copyright see the source file.
%% 
%% Any modified versions of this file must be renamed
%% with new filenames distinct from sample-manuscript.tex.
%% 
%% For distribution of the original source see the terms
%% for copying and modification in the file samples.dtx.
%% 
%% This generated file may be distributed as long as the
%% original source files, as listed above, are part of the
%% same distribution. (The sources need not necessarily be
%% in the same archive or directory.)
%%
%% The first command in your LaTeX source must be the \documentclass command.
%%%% Small single column format, used for CIE, CSUR, DTRAP, JACM, JDIQ, JEA, JERIC, JETC, PACMCGIT, TAAS, TACCESS, TACO, TALG, TALLIP (formerly TALIP), TCPS, TDSCI, TEAC, TECS, TELO, THRI, TIIS, TIOT, TISSEC, TIST, TKDD, TMIS, TOCE, TOCHI, TOCL, TOCS, TOCT, TODAES, TODS, TOIS, TOIT, TOMACS, TOMM (formerly TOMCCAP), TOMPECS, TOMS, TOPC, TOPLAS, TOPS, TOS, TOSEM, TOSN, TQC, TRETS, TSAS, TSC, TSLP, TWEB.
% \documentclass[acmsmall]{acmart}

%%%% Large single column format, used for IMWUT, JOCCH, PACMPL, POMACS, TAP, PACMHCI
% \documentclass[acmlarge,screen]{acmart}

%%%% Large double column format, used for TOG
% \documentclass[acmtog, authorversion]{acmart}

%%%% Generic manuscript mode, required for submission
%%%% and peer review
%\documentclass[manuscript,screen,review]{acmart}
\documentclass[conference]{IEEEtran}
%\documentclass[sigconf,review]{acmart}


%% Fonts used in the template cannot be substituted; margin 
%% adjustments are not allowed.
%%
%% \BibTeX command to typeset BibTeX logo in the docs
\AtBeginDocument{%
 \providecommand\BibTeX{{%
   \normalfont B\kern-0.5em{\scshape i\kern-0.25em b}\kern-0.8em\TeX}}}
    
  %  \AtBeginDocument{\renewcommand*{\thesubfigure}{\alphalph{\value{subfigure}}}}

%% Rights management information.  This information is sent to you
%% when you complete the rights form.  These commands have SAMPLE
%% values in them; it is your responsibility as an author to replace
%% the commands and values with those provided to you when you
%% complete the rights form.


% \setcopyright{none} % remove permission section
% \settopmatter{printacmref=false} % remove ACM reference format
% \renewcommand\footnotetextcopyrightpermission[1]{} % removes footnote with conference information in first column
\pagestyle{plain}
%\setcopyright{acmcopyright}
%\copyrightyear{2022}
%\acmYear{2022}
%\acmDOI{00.0000/0000000.0000000}

%% These commands are for a PROCEEDINGS abstract or paper.
%\acmConference[Woodstock '18]{Woodstock '18: ACM Symposium on Neural Gaze Detection}{June 03--05, 2018}{Woodstock, NY}
%   \acmConference[MSR 2022]{MSR '22: Proceedings of the 19th International Conference on Mining Software Repositories}{May 23–24, 2022}{Pittsburgh, PA, USA}
% \acmBooktitle{MSR '22: Proceedings of the 19th International Conference on Mining Software Repositories, May 23–24, 2022, Pittsburgh, PA, USA}
% \acmPrice{15.00}
% \acmISBN{978-1-4503-XXXX-X/18/06}

\usepackage{xspace}
\newcommand{\ie}{\textit{i.e.}}
\newcommand{\eg}{\textit{e.g.}}
\newcommand{\etal}{\textit{et al. \xspace}}
%\usepackage{array}
%\newcolumntype{L}{>{\arraybackslash}m{8cm}}
%%%%% Long text wrap %%%%%
\usepackage{array}
\newcolumntype{L}{>{\arraybackslash}m{16cm}}
%\newcolumntype{L}[1]{>{\raggedright\let\newline\\arraybackslash\hspace{0pt}}m{#1}}
\newcolumntype{C}[1]{>{\centering\let\newline\\arraybackslash\hspace{0pt}}m{#1}}
\newcolumntype{R}[1]{>{\raggedleft\let\newline\\arraybackslash\hspace{0pt}}m{#1}}
%%%%% Long text wrap %%%%%
%\usepackage{adjustbox}
%\usepackage{subcaption}
%\usepackage{textcomp}
%\usepackage{pgfplots}
\usepackage[numbered]{bookmark}
%\usepackage{soul}
%\usepackage{booktabs}
%\usepackage{multirow}
%\usepackage{float}
%\usepackage{url}
\usepackage{tcolorbox}
\usepackage{graphicx}
%\usepackage{textcomp}
\usepackage{xcolor}
%\usepackage{url}
%\usepackage{tabularx} 
%\usepackage{balance}
\usepackage{stfloats}
%\usepackage{fullpage}
\usepackage[justification=centering]{caption}
\usepackage{lineno}
% \usepackage[title]{appendix}
\usepackage{dirtytalk}
\usepackage{amsmath}
%\usepackage{tikz}
\usepackage{pgfplots, pgfplotstable}
%\usetikzlibrary{tikzmark}
\usetikzlibrary{patterns}



%%%%%
%-----------------------
%\usepackage[utf8]{inputenc}
\usepackage[english]{babel}
%\usepackage{amssymb}
\usepackage{textcomp}
%\usepackage{indentfirst}
%\setcounter{tocdepth}{3}
\usepackage{mathtools}
%\usepackage[usestackEOL]{stackengine}
%\usepackage{amsmath}
\usepackage{ltxtable}
\usepackage{algorithm}
\usepackage{algpseudocode}
\usepackage{textcomp}
%\usepackage{mathpazo}
%\usepackage{tgpagella}
%-----------------------

\usepackage{tcolorbox}

%\usepackage[dvipsnames]{xcolor}
\usepackage{pgf-pie}
%\usepackage{pgfplots}
%\pgfplotsset{compat=1.14}
%%%%%%%%%%%%%%%%%%%%%%%%%%%%%%%%%%%%%%%%%%%%%%%%%%%
\definecolor{wedge1}{RGB}{ 190  30  46}
\definecolor{wedge2}{RGB}{ 240  65  54}
\definecolor{wedge3}{RGB}{ 241  90  43}
\definecolor{wedge4}{RGB}{ 247 148  30}
\definecolor{wedge5}{RGB}{  43  56 144}
\definecolor{wedge6}{RGB}{  28 117 188}
\definecolor{wedge7}{RGB}{  40 170 225}
\definecolor{wedge8}{RGB}{ 119 179 225}
\definecolor{wedge9}{RGB}{ 181 212 239}
\definecolor{wedge10}{RGB}{  0 104  56}
\definecolor{wedge11}{RGB}{  0 148  69}
\definecolor{wedge12}{RGB}{ 57 181  74}
\definecolor{wedge13}{RGB}{141 199  63}
\definecolor{wedge14}{RGB}{215 244  34}
\definecolor{wedge15}{RGB}{249 237  50}
\definecolor{wedge16}{RGB}{248 241 148}
\definecolor{wedge17}{RGB}{242 245 205}
\definecolor{wedge18}{RGB}{123  82  49}
\definecolor{wedge19}{RGB}{104  73 158}
\definecolor{wedge20}{RGB}{102  45 145}
\definecolor{wedge21}{RGB}{148 149 151}
\definecolor{wedge22}{RGB}{ 204 50 153}
\definecolor{wedge23}{RGB}{ 79 47 79}
\definecolor{wedge24}{RGB}{ 173 234 234}
\definecolor{wedge25}{RGB}{ 216 191 216}
\definecolor{wedge26}{RGB}{  43  56 144}
\definecolor{wedge27}{RGB}{  40 170 225}
\definecolor{wedge28}{RGB}{ 119 179 225}
\definecolor{wedge29}{RGB}{ 181 212 239}
\definecolor{wedge30}{RGB}{  0 104  56}
\definecolor{wedge31}{RGB}{  0 148  69}
\definecolor{wedge32}{RGB}{ 57 181  74}

\usetikzlibrary{chains}
%%%%%%%%%%%%%%%%%%%%%%%%%%%%%%%%%%%%%%


\def\groupTotals{{1,4,5,8,1,2}}
\pgfmathsetmacro\startAngle{90-3.6/2}
\pgfmathsetmacro\radius{+5}
\pgfmathsetmacro\maxLeg{+12}
\pgfmathsetmacro\legBound{+60}
\pgfmathsetmacro\legSpacing{2*\legBound/(\maxLeg-1)}
%%%%%%%%%%%%%%%%%%%%%%%%%%%%%%%%%%%%%%

%%%%%%%%%%%%%%%%%%%%%%%%%%%%%%%%%%%%%
\usepackage{verbatim}
%\usepackage{pgfplots, pgfplotstable}
%\pgfplotsset{width=7cm,compat=1.8}
%\pgfplotsset{width=7cm,compat=1.8,tick label style={font=\small}}
\pgfplotsset{compat=1.14}
\usetikzlibrary{patterns,}

\usepackage{graphicx}
\usepackage{tabularx,booktabs}
\usepackage{array}
\usepackage{bbding}
\usepackage{pifont}
\usepackage{wasysym}
%\usepackage{amssymb}
\usepackage{caption}
\usepackage{dirtytalk}
%\usepackage[showframe]{geometry}
%\usepackage{subfigure}
\usepackage{subcaption}
%\usepackage{amsmath} % Star
%\usepackage[table]{xcolor}
\usepackage{tabularx}
%\usepackage{pgfplots}
%\usepackage{pgf-pie}
%\usepackage{verbatim}
%\usepackage{pgfplotstable}
\usepackage{sfmath}
\usepackage{array, booktabs, makecell}
\usepackage{color}
%\usepackage[dvipsnames]{xcolor}
\usepackage{csquotes}
\usepackage{url}
\usepackage{comment}
\usepackage{tikz}
\usepackage{balance}
\usepackage{listings}
\usepackage{booktabs} % used for \toprule in tables
%\usepackage{cite}
\usepackage{soul} % highlighting
\urlstyle{same} %% footnote fonts
\usetikzlibrary{fit} %% Used for putting dotted box in image
%\usepackage[margin=2cm]{geometry}
\usetikzlibrary{positioning}
\usetikzlibrary{arrows}
\usetikzlibrary{shapes.multipart}
%\usepackage{caption}
%\usepackage{graphicx}
\usepackage{adjustbox}
%\usepackage{tablefootnote}
\usepackage{float}
\usepackage{rotating}
\usepackage{nth}
\DeclareCaptionType{TextBox}
\usepackage{pstricks}
\usepackage{alltt}
\usepackage{microtype}
%%%%%%%%%%%%%%%%%%%%%%%%%%%%%%%%%%%%%%%%%%%%%%%%%%%%%%%%
\usepackage{placeins}
\usepackage{afterpage}
\usepackage{multirow} 
\usepackage{textcomp}
\usepackage{multicol}
\usepackage{varwidth} % Double column tables

%%%%%

%\usepackage{subfigure}
\newcommand{\eman}[1]{\textcolor{violet}{{\it [Eman says: #1]}}}




%%
%% Submission ID.
%% Use this when submitting an article to a sponsored event. You'll
%% receive a unique submission ID from the organizers
%% of the event, and this ID should be used as the parameter to this command.
%%\acmSubmissionID{123-A56-BU3}

%%
%% The majority of ACM publications use numbered citations and
%% references.  The command \citestyle{authoryear} switches to the
%% "author year" style.
%%
%% If you are preparing content for an event
%% sponsored by ACM SIGGRAPH, you must use the "author year" style of
%% citations and references.
%% Uncommenting
%% the next command will enable that style.
%%\citestyle{acmauthoryear}

%%
%% end of the preamble, start of the body of the document source.
\begin{document}

%%
%% The "title" command has an optional parameter,
%% allowing the author to define a "short title" to be used in page headers.
\title{\huge Beyond Test Presence: Assessing the Quality and Robustness of Agent-Generated Tests in Open-Source Projects}

 \author{
   \IEEEauthorblockN{Preet Jhanglani\IEEEauthorrefmark{1},
                     Zeel Kaushal Desai\IEEEauthorrefmark{1},
                     Vidhi Kansara\IEEEauthorrefmark{1},
                     Eman Abdullah AlOmar}
   \IEEEauthorblockA{Stevens Institute of Technology, 
                     Hoboken, New Jersey, USA\\
                     Email: \{pjhangl1,zdesai,vkansara,ealomar\}@stevens.edu}
 }



% \author{Preet Jhanglani}
% %\authornote{Both authors contributed equally to this research.}
% \email{pjhangl1@stevens.edu}
% \affiliation{%
%   \institution{Stevens Institute of Technology}
%   %\streetaddress{P.O. Box 1212}
%   \city{Hoboken}
%   \state{New Jersey}
%   \country{USA}
%   %\postcode{43017-6221}
% }


  
% \author{Zeel Kaushal Desai}
% %\authornote{Both authors contributed equally to this research.}
% \email{zdesai@stevens.edu}
% \affiliation{%
%   \institution{Stevens Institute of Technology}
%   %\streetaddress{P.O. Box 1212}
%   \city{Hoboken}
%   \state{New Jersey}
%   \country{USA}}
  
% \author{Vidhi Kansara}
% %\authornote{Both authors contributed equally to this research.}
% \email{vkansara@stevens.edu}
% \affiliation{%
%   \institution{Stevens Institute of Technology}
%   %\streetaddress{P.O. Box 1212}
%   \city{Hoboken}
%   \state{New Jersey}
%   \country{USA}}
  

%%
%% The "author" command and its associated commands are used to define
%% the authors and their affiliations.
%% Of note is the shared affiliation of the first two authors, and the
%% "authornote" and "authornotemark" commands
%% used to denote shared contribution to the research.


%%
%% By default, the full list of authors will be used in the page
%% headers. Often, this list is too long, and will overlap
%% other information printed in the page headers. This command allows
%% the author to define a more concise list
%% of authors' names for this purpose.
%\renewcommand{\shortauthors}{Trovato and Tobin, et al.}

%%
%% The abstract is a short summary of the work to be presented in the
%% article.
\maketitle
\begingroup
\renewcommand\thefootnote{\IEEEauthorrefmark{1}}
\footnotetext{These authors contributed equally to this work.}
\endgroup
\begin{abstract}
The integration of AI-powered coding agents into Continuous Integration/Continuous Delivery (CI/CD) pipelines has fundamentally altered how software verification is conducted. While these agents successfully automate the test generation, current evaluation benchmarks (e.g., SWE-bench) largely focus on pass-rates rather than the intrinsic quality of the generated tests. This raises the possibility of ``stealth technical debt'', in which test suites pass execution but do not offer comprehensive coverage or semantic value. We address this methodological gap through a large-scale, empirical comparison of 204,673 test artifacts which comprises of 24,941 human-authored files and 179,732 agent-generated files; sourced from the AIDev dataset.

Using the Abstract Syntax Tree (AST) parsing with Python's naive \texttt{ast} module, we implemented a ``white-box'' static analysis framework to evaluate three quality dimensions: Assertion Strength (RQ1), Edge-Case Coverage (RQ2), and Flakiness Potential (RQ3). Our results present a nuanced inversion of traditional assumptions. AI agents performed better than humans in Edge-Case Coverage, with almost twice the variety of boundary checks (Variety Score: 0.62 vs 0.32) and a higher frequency of null-safety testing (13.40\% vs. 8.3\%), even though human developers had a slight advantage in Assertion Strength (88.1\% strong assertions vs. 85.37\% for agents). But this thoroughness comes at a price: due mostly to their reliance on file I/O and non-deterministic logic, agent-generated tests exhibited a higher risk of flakiness (Candidate Rate: 0.41 vs. 0.30). These findings suggest that while AI agents excel at rigorous boundary testing, they lack the ``environmental awareness'' needed to write stable, hermetic tests.
\end{abstract}
%%
%% The code below is generated by the tool at http://dl.acm.org/ccs.cfm.
%% Please copy and paste the code instead of the example below.
%%
% \begin{CCSXML}
% <ccs2012>
%  <concept>
%   <concept_id>10010520.10010553.10010562</concept_id>
%   <concept_desc>Computer systems organization~Embedded systems</concept_desc>
%   <concept_significance>500</concept_significance>
%  </concept>
%  <concept>
%   <concept_id>10010520.10010575.10010755</concept_id>
%   <concept_desc>Computer systems organization~Redundancy</concept_desc>
%   <concept_significance>300</concept_significance>
%  </concept>
%  <concept>
%   <concept_id>10010520.10010553.10010554</concept_id>
%   <concept_desc>Computer systems organization~Robotics</concept_desc>
%   <concept_significance>100</concept_significance>
%  </concept>
%  <concept>
%   <concept_id>10003033.10003083.10003095</concept_id>
%   <concept_desc>Networks~Network reliability</concept_desc>
%   <concept_significance>100</concept_significance>
%  </concept>
% </ccs2012>
% \end{CCSXML}

% \ccsdesc[500]{Software Engineering~Software Testing, Quality Assurance, and Maintenance}
% \ccsdesc[300]{Software Engineering}
%% \ccsdesc{Computer systems organization~Robotics}
%% \ccsdesc[100]{Networks~Network reliability}

%%
%% Keywords. The author(s) should pick words that accurately describe
%% the work being presented. Separate the keywords with commas.
%\keywords{software quality, mining software repositories}
\begin{IEEEkeywords}
testing, quality, agents, pull requests
\end{IEEEkeywords}
%% A "teaser" image appears between the author and affiliation
%% information and the body of the document, and typically spans the
%% page.

%%
%% This command processes the author and affiliation and title
%% information and builds the first part of the formatted document.


\section{Introduction}
\label{Section:Introduction}
At the present time, technology is growing so quickly that almost every application receives continuous feature updates. Companies are now releasing regular updates at a much faster pace than before, and many times a day. Once a feature or product has been developed, that feature or product will go through a period of testing to ensure that it passes all test cases prior to actual release. A majority of companies take advantage of continuous integration and continuous deployment (CI/CD) pipelines, which allow for the automation of software development and testing. CI/CD pipelines can be integrated with sites like GitHub, where each pull request is automatically generated and tested before being incorporated into the main code structure.

%%At the present time, technology is growing very fast, and most applications keep getting new features regularly. Compared to earlier times, companies now release updates much more frequently. Sometimes updates happen daily. After a feature or product is developed, it goes for testing. If it passes all the test cases, then it is released. CI/CD pipelines are used by many businesses to automate software development and testing. These pipelines are compatible with sites such as GitHub, where each pull request is automatically produced and tested prior to being pushed to the main code branch.

 Testing is a key part of the development process, even for small updates. The developer needs to determine how the code should behave and identify the most likely situations that could result in breaking the code. The developer should keep these test cases current as the codebase continues to expand. In organizations that release frequently, there may be insufficient time to write adequate test cases and ensure that tests are executed correctly. Among other reasons, one of the main reasons why developers are beginning to leverage AI-powered coding tools is to automate sections of the development and testing life cycle, which also allows them to have more available time to write tests. Developers can use AI-powered coding tools, such as GitHub Copilot and more autonomous agents like Devin, to read a project description, generate implementation code, and also auto-generate test cases to go along with that code. Essentially, they act as junior members of a team that develops a feature as well as creates the initial test coverage, thus reducing the amount of effort developers would typically spend doing work manually \cite{yoshimoto2026testing,yuan2023no}.

%%Testing is a very important part of development, even for small updates. Developers need to think through how the code is supposed to behave, identify the cases most likely to break it, and keep those tests up to date as the codebase grows. In teams that release frequently, finding time for careful test authoring can be a real challenge. This is one reason AI-powered coding tools have started playing a bigger role. Tools like GitHub Copilot, and more autonomous agents such as Devin, can read a task description, write implementation code, and generate matching test cases alongside it. In practice, they work like junior team members who handle both the feature and its initial test coverage, reducing the manual effort developers would otherwise spend \cite{compsac_replication_package}.

There is a growing interest in researching the application of AI to generate test cases \cite{siddiq2024using,gao2025current,celik2025review,ouedraogo2024large,deljouyi2024leveraging,molina2025test,pan2024multi,qiu2025today,milanese2026human}; for example, one well-known project is TestPilot, which uses an algorithmatically generated JavaScript unit test from a function's source code, signature, and usage examples by feeding the code into a large language model (LLM) to create the test case automatically. After being tested with 25 NPM packages, the median line coverage for the generated unit tests was found to be 70.2\%, and the median branch coverage was 52.8\%. This provides a meaningful increase over previous automated testing solutions for generating unit test cases. Subsequent studies such as CoverUp have provided further refinements including the ability to identify which segments of code have not been unit tested yet and have the algorithm concentrate on generating test cases for those areas \cite{altmayer2025coverup}. The majority of test cases generated by this approach have been evaluated using benchmark suites such as SWE-Bench \cite{jimenez2023swe} to determine whether the generated patches pass developer-created unit test suites. However, there are studies indicating that many generated patches that were indicated as passing on their behalf were still lacking in terms of the original intent and/or were not in line with the expected behavior. This indicates that while all tests were indeed successful, the resulting patch could still have issues and/or could not sufficiently resolve the issue originally intended.

%%Researchers have studied AI-driven test generation with growing interest. One notable example is TestPilot, which uses a large language model (LLM) to produce JavaScript unit tests by feeding it a function's code, signature, and usage examples. Tested against 25 npm packages, it achieved a median statement coverage of 70.2\% and branch coverage of 52.8\%, which is meaningfully higher than older automated tools. Later work, such as CoverUp, improved on this by identifying which parts of the code were not yet covered and asking the model to specifically target those gaps [1]. Alongside these tools, benchmarks such as SWE-bench are commonly used to evaluate whether AI-generated patches can pass test suites written by developers. However, studies like UTBoost have shown that many patches marked as “passing” are still incomplete or semantically incorrect. This shows that even if every test is successful, the resulting patch might still have issues or might not resolve the problem completely \cite{compsac_replication_package}.

The quality of test cases can be affected by several other factors, such as assertion strength. Assertion strength relates to how stringent or lenient a test is in its verification of the result. Here is one example of a good assertion strength -- comparing the exact result produced by the program as opposed to only verifying whether true was returned. Edge-case coverage also plays a significant role in the quality of the test case itself. Good tests should contain edge-case coverage for abnormal input values (e.g., nulls, empty lists, 0, max val, min val) since some of the most elusive defects are located in these types of input cases. Test case stability is the third area impacting test case quality; stability refers to how consistently a test produces the same outcome when executed repeatedly. Unpredictable outcomes (i.e., both passing and failing) caused by flaky tests destroy the confidence level of the test suite and increase the time required to release. During periods of fast-paced development, the quality of these three areas, regardless of whether the test case was manually written or created by a test generation tool, tend to decline \cite{barraood2021test,crispin2006driving}.

%Some other things also affect the quality of test cases. One important thing is assertion strength. This means how carefully the test checks the result. For example, checking the exact output is better than only checking whether the program returns `true`. Another important thing is edge-case coverage. Good tests should also check unusual inputs such as `null` values, empty lists, zero, or maximum and minimum values, because these cases can often find hidden bugs in software. The third is test stability, or whether the test gives the same result every time it runs. Tests that pass and fail unpredictably, often called flaky tests, are a serious problem because they erode trust in the test suite and slow down the release process. When developers are working quickly, all three of these qualities can suffer, whether the tests are written by a person or generated by a tool \cite{compsac_replication_package}.

There is an important consideration when it comes to AI-generated tests. Most comparisons of AI-generated testing focus on the percentage of lines of code that the tests cover, or whether the tests were successful or not. These are both reasonable baseline comparisons, but they add very little information about how reliable the tests actually are. A test suite can have excellent coverage, but may have missed some important edge cases, contain weak assertions, or contain unreliable tests that may fail unpredictably. If testing tools are only designed to evaluate based on coverage measurements, they will produce tests that seem reliable but provide very little protection once the software is in production \cite{yang2006survey}.

By concentrating on assessing the quality of tests directly in relation to all three of the above measurements; human written versus AI generated, through the use of actual real-world software development activities contained in the AIDev database (the largest database currently available of actual software development activity). The tests are not actually executed; rather, abstract syntax trees (AST) will be parsed from each test so that structure can be evaluated statically and thus be able to scale and provide consistent results. There are 179,732 AI-generated tests (files) and 24,941 human-created tests (files) present in this database. Each of these files contained in the AIDev database went through three independent automated test evaluation pipelines in our evaluation process.

The findings provide a mixed result in terms of test strength. For example, 88 of the human-written assertions were categorized as strong; however, the total number of ambiguous or unclassified assertions recorded by human testers was far less than those written by computers. Human-written tests were able to use only 1.46\% ambiguous or unclassified assertions while AI-generated tests had 11.58\%. On edge-case coverage, the situation is reversed. AI agents covered a broader variety of edge cases, scoring 0.62 on the variety measure against 0.32 for humans. Humans performed better than AI – Agents on test stability, this study showing humans generated tests achieved a flakiness rate of 0.30 compared to 0.41 for AI-generated tests.

For this study, we worked on these three research questions:
\begin{enumerate}
    \item \textbf{RQ1: In comparison to tests created by human developers, how solid are the claims in test cases created by AI?}
    \item \textbf{RQ2: How well can AI-generated test cases manage boundary conditions and edge situations in comparison to tests written by humans?}
    \item \textbf{RQ3: Do flaky tests appear more frequently in pull requests made by human developers or by AI tools?}
\end{enumerate}
The following sections of this work are organized as follows. Section \ref{Section:Methodology} describes the approach, as well as the procedures used to collect and analyze the data. Section \ref{Section:Result} discusses the results for each study question. Potential validity risks are described in Section \ref{Section:Threats}. The conclusion and future work are provided in Sections \ref{Section:Conclusion} and \ref{Section:FutureWork}.

% The world of software development is changing rapidly these days. New features, enhancements, and bug fixes are planned for release by teams in brief cycles, sometimes within days or even hours. This can be achieved through continuous integration and continuous delivery (CI/CD) pipelines. These systems automatically generate and test software whenever new code is added, prior to its release. Automated tests, such as unit, integration, and regression tests, play an important role in this context. They help catch errors early and prevent small issues from becoming major problems.

% Although testing is very important, writing tests takes time and effort. Developers have to not only write new features, but also think about edge cases, error situations, and unusual inputs. Writing tests and code at the same time may slow development. Because of this, AI-powered coding tools are becoming more common. Modern AI agents can do more than just suggest code. They can understand a task, break it into steps, generate code, run checks, and even create pull requests. They behave much like junior developers who help human engineers.

% Automated test generation is a primary application of these AI agents. Many teams request that new code be tested before it can be merged. It helps preserve test coverage and saves time if an AI agent can create tests while a feature is being implemented. The goal is for testing to become part of the development cycle rather than an isolated activity.

%In modern software engineering, change is the only constant.  Development teams are constantly under pressure to produce new features, improvements, and fixes in ever-faster cycles; often in a matter of hours or days as opposed to weeks. The well-oiled machinery of CI/CD pipelines is responsible for this acceleration. Every time new code is committed, these pipelines make sure it is automatically built and tested; it does not go to deployment until it passes the validation suite. Automated tests including unit, integration and regression serve as crucial gatekeepers in such a setting.They detect regressions and unexpected interactions early, before those issues proliferate into larger, more costly failures. However, a paradox still exists.  Even though tests are essential, many people find writing them to be time-consuming, particularly when every new or changed piece of code needs to have matching tests. Developers must list edge cases, error routes, and unusual interactions in addition to reasoning through the proper behavior of new code. In fast-paced workflows, the mental shift between feature development and test authoring can cause friction and reduce productivity. Coding bots driven by AI are increasingly filling this role. More sophisticated autonomous agents like Devin and more recent iterations of programs like GitHub Copilot go beyond straightforward code completions.  They are able to take a task description, break it down into smaller sections, create implementation code, perform initial checks, and even initiate pull requests for review. In effect, they begin to behave like junior collaborators embedded within the CI/CD pipeline, bridging the gap between human intention and executable patches. One of the most compelling roles for such agents is automated test generation. In many organizations, reviewers insist that new or changed code arrives bundled with tests. If an AI agent can generate plausible tests concurrently with feature implementation, it reduces developer overhead and helps maintain or improve coverage. In principle, test writing is no longer a separate chore; it becomes part of the development flow, meaning every new capability arrives with at least minimal protection.

% However, this convenience introduces a subtle but serious tension: how do we ensure that speed does not undermine quality? The promise of CI/CD is fast, safe shipping but that promise depends on the discriminative power of the the test suite guarding each commit. If AI-generated tests are superficial; covering only normal, easy paths and failing to account for boundary conditions, invalid inputs, or failure modes then a continuously ``green'' build might become misleading. Over time, a codebase can accumulate latent defects that evade detection, creating what practitioners call ``false security'' or ``stealth technical debt.'' 

% While large-scale datasets like AIDev provide an invaluable starting point for studying developer behavior, they often capture development artifacts as static snapshots or differential patches[11]. To enable deep semantic analysis of source code such as the Abstract Syntax Tree (AST) parsing required for our work. It is essential to reconstruct the complete state of software artifacts at precise points in their history. We therefore employ a two-phase data collection strategy: first identifying key development events from the AIDev corpus, then programmatically retrieving the full content of associated test files from live source repositories to ensure the fidelity of our subsequent analysis.

%To understand the current capabilities and limitations of AI-driven test generation, researchers have been studying large language model (LLM) approaches in empirical settings \cite{siddiq2024using,gao2025current,celik2025review,ouedraogo2024large,deljouyi2024leveraging,molina2025test,pan2024multi}. A prominent example is TestPilot, which prompts an LLM (e.g., GPT-3.5) with a function's implementation, signatures, and usage snippets to synthesize JavaScript unit tests. In an evaluation of 25 npm packages (1,684 API functions), the authors report a median statement coverage of 70.2\% and branch coverage of 52.8\%, significantly exceeding an older tool, Nessie \cite{yang2024empirical}. They also adopt a feedback loop: if a generated test fails (e.g., does not compile or triggers an error), the agent is re-prompted with failure information to repair or refine the test. Subsequent techniques have extended this work. For instance, CoverUp introduces a coverage-guided prompting strategy, where the system identifies under-covered regions and explicitly asks the model to target those areas in subsequent prompts, significantly increasing coverage over naive test generation.Many such methods combine prompt engineering, iterative repair, and heuristics to drive models toward greater coverage and fewer errors.

%To benchmark the broader capabilities of AI coding agents, many works adopt SWE-bench, a benchmark that tasks models with resolving real GitHub issues by generating patches that pass existing test suites. Although widely used, SWE-bench \cite{jimenez2023swe} mainly measures whether a patch passes human-written tests; it does not systematically evaluate the robustness or fault-detection capability of newly generated tests. Recent work, such as UTBoost \cite{yu2025utboost}, has revealed that many patches accepted by the original SWE-bench test suites are semantically incorrect or incomplete. UTBoost uncovered 345 erroneous patches wrongly marked as passing, affecting the rankings of SWE-bench Lite and Verified \cite{jimenez2023swe}. These findings suggest that even benchmark-based evaluations can understate the fragility of test suites constructed without deeper scrutiny.

%In sum, we observe a methodological gap. AI test generation is gaining traction and improving in coverage, but evaluation practices largely emphasize coverage and pass/fail outcomes; metrics that are easy to compute but often poor proxies for true test quality. If left unchecked, this bias may push AI systems to optimize for superficial metrics rather than semantic depth, edge-case coverage, or real fault detection.The result could be a generation of test suites that look polished in dashboards but provide weak protection in practice. 

%In this paper, we seek to address this gap by proposing an evaluation framework that emphasizes fault-finding capability, assertion richness, and semantic robustness of generated tests, in addition to conventional metrics. We apply it to test suites produced by AI agents in real-world open-source pull requests, compare how different evaluation lenses reveal hidden weaknesses, and offer insights for designing more reliable test generation systems and benchmarks. Specifically, we investigate three research questions (RQs):
% \section{Problem Statement}
% \label{Section:Problem}
% The proliferation of AI-powered coding agents, from code assistants to autonomous developers, has fundamentally altered the software engineering landscape. While a large body of research has focused on the functional correctness of agent-generated code, a methodological gap persists. Benchmarks like SWE-bench primarily measure whether an AI-generated patch causes an existing, human-written test suite to pass \cite{Bellen2013}.This evaluation paradigm overlooks a critical question: what is the quality of the tests themselves when generated by AI?
% If AI agents, in their quest to satisfy CI/CD pipelines, produce tests that are superficial, brittle, or non-deterministic, they may introduce "stealth technical debt" and create "false security" [1]. This study addresses this gap by directly comparing the intrinsic quality of AI-generated and human-authored test artifacts. 

% We seek to answer this through an evaluation framework that emphasizes fault-finding capability, assertion richness, and semantic robustness of generated tests, in addition to conventional metrics.
% This research is architected as an extensive, large-scale comparative analysis, designed to meticulously scrutinize the qualitative and quantitative differences between software tests authored by human developers and those generated by autonomous AI agents. 


%\begin{itemize}
        %\item \textbf{RQ1: How does the assertion strength of agent-generated tests compare to human-written tests?}
        %: Assertion strength reflects the semantic depth of validation checks. We hypothesize that human developers, drawing on domain expertise, will employ more precise assertions that verify specific behavioral invariants, while agents may default to shallow existence checks.
        %\item \textbf{RQ2: To what extent do agent-generated tests cover edge and boundary cases compared to human-written tests?}
        %: Edge-case coverage is a critical indicator of test thoroughness. We seek to determine whether agents systematically enumerate boundary conditions (null, zero, empty collections) or whether they focus primarily on "happy path" scenarios.
       % \item \textbf{RQ3: What is the prevalence of flaky tests in pull requests submitted by AI agents compared to those submitted by humans?} 
        %: Test flakiness and non-deterministic test outcomes undermines the CI/CD reliability. We investigate whether agent-generated tests exhibit higher rates of flakiness-inducing patterns such as unmanaged randomness, fixed time delays, or environmental dependencies.
%\end{itemize}

% \section{Related Work}
% \label{Section:RelatedWork}
% The proliferation of AI-powered coding agents has fundamentally altered the software engineering landscape. While much research has focused on the functional correctness of agent-generated code, a methodological gap persists: benchmarks like SWE-bench primarily measure whether an AI-generated patch causes an existing, human-written test suite to pass\cite{Bellen2013}], but this evaluation paradigm overlooks a critical question—what is the quality of the tests themselves when generated by AI? If AI agents, in their quest to satisfy CI/CD pipelines, produce tests that are superficial, brittle, or non-deterministic, they may introduce "stealth technical debt" and create "false security" [1]. We organize the related literature into three key areas that inform our study design.

% \subsection{Empirical Studies of LLM vs. Human Code Quality}

% Recent (2024–2025) research has begun to shift from "can AI code?" to "is AI-generated code good?" This line of inquiry provides a direct foundation for our project. A key study by Licorish et al. (2025) \cite{Licorish2025} performs a comparison of human- and LLM-generated code using metrics of functional correctness, code quality, and complexity. A significant finding is that LLM-generated code, while often functionally correct, "frequently produces more complex code" (e.g., higher cyclomatic complexity) that may "need more reworking to ensure maintainability". This finding justifies our need to analyze the maintainability and quality of AI-generated tests, as higher complexity in implementation code may be mirrored by higher complexity or "smells" in test code.

% Further supporting this is the work of Molison et al. (2025) [8], "Is LLM-Generated Code More Maintainable and Reliable than Human-Written Code?". This study provides a critical, nuanced conclusion: LLM-generated code was found to have "fewer bugs" and "require less effort to fix them" in simple scenarios, but it also introduced "critical issues in more complex scenarios." This finding directly aligns with our hypotheses: AI-generated tests may excel at the "happy path" (fewer simple bugs) but fail to cover complex edge cases (RQ2) or may naively introduce "critical issues" like resource leaks or non-determinism—root causes of flakiness (RQ3). 

% \subsection{Studies on Automated Test Generation}
% The automated generation of unit tests is a long-standing field in software engineering. Traditional, non-LLM approaches such as search-based (e.g., EvoSuite) or symbolic-execution-based tools have been widely studied. However, this body of work has consistently found that while such tools can achieve high coverage, they "still fall short of the high utility achieved by human-written tests," often attributed to the generation of tests with "code smell issues" that are "difficult to maintain and understand."

% The emergence of LLMs has been seen as a promising solution to this "low utility" problem. Yang et al. (2024)\cite{Yang2024} conducted "An Empirical Study of Unit Test Generation with Large Language Models," one of the first to compare open-source LLMs, GPT-4, and traditional tools like EvoSuite. Their findings highlight that LLM performance is highly sensitive to "various prompt factors" and that, despite their promise, "limitations in LLM-based unit test generation" remain. Our project builds directly on this work, moving beyond coverage and pass-rate metrics to evaluate the qualitative limitations—such as weak assertions, poor edge-case coverage, and flakiness—that may be present in LLM-generated tests.

%\subsection{Methodologies for Test Quality Analysis}
%Our methodology, which uses AST-based static analysis to quantify test quality, is grounded in established test quality analysis techniques that justify each of our three research questions.

    %\begin{itemize}
       % \item \textbf{RQ1 (Assertion Strength):}
        %The premise of RQ1—that assertion strength is a meaningful proxy for test quality—is empirically supported. The foundational work by Zhane and Mesbah  \cite{zhang2015assertions} investigated the relationship between test assertions and test suite effectiveness (measured by mutation score).The study found that assertion quantity and type are "strongly correlated" with a test suite's effectiveness. This directly validates our focus on assertion analysis. Critically, that study also provided a counter-intuitive finding: \texttt{assertTrue/False assertions} were, on average, more effective then \texttt{assertEquals/Not assertions}. This is because \texttt{assertTrue} is often used to encapsulate complex domain logic, whereas \texttt{assertEquals} is often used for simple state checks. This finding directly informs the ``Advanced Taxonomy'' we use in our white-box design. %(see Section 4).
       % \item \textbf{RQ2 (Edge-Case Coverage):} Edge-case coverage is a critical indicator of test thoroughness. We seek to determine whether agents systematically enumerate boundary conditions (null, zero, empty collections). 
       % \item \textbf{RQ3 (Test Flakiness):}
       % The methodology for RQ3, which uses static analysis to find ``flakiness candidates,'' is a well-known technique in test smell and flakiness detection research. Numerous studies (e.g., \cite{peruma2020tsdetect,bavota2015test,gruber2024automatic}) have demonstrated the use of AST parsing to find test anti-patterns or ``test smells.'' Other work has successfully used static metrics and test smells as predictive features for machine learning models to identify flaky tests \cite{dutta2020detecting}. Other research identifies ``Async Wait, Concurrency, Test Order Dependency, and Resource Leak'' as key root causes of flakiness that can serve as predictive features. Therefore, the static analysis in Pipeline C represents a proven, valid method for identifying tests that have a high statistical likelihood of being flaky.
   % \end{itemize}

%The remainder of this paper is organized as follows:  Section \ref{Section:Methodology} outlines our empirical setup
% in terms of data collection, analysis, and research question.
%Section \ref{Section:Result} discusses our findings, while Section \ref{Section:Threats} captures
%any threats to the validity of our work, before concluding with
%Sections \ref{Section:Conclusion} and \ref{Section:FutureWork}.

% \begin{table*}[h!]
% \centering
% \caption{Summary of Key Related Work}
% \begin{tabular}{|p{3.5cm}|p{3.5cm}|p{4cm}|p{5cm}|}
% \hline
% \textbf{Study} &
% \textbf{Methodology} &
% \textbf{Context} &
% \textbf{Key Findings Relevant to This Project} \\
% \hline

% Licorish et al. 2025 \cite{Licorish2025} &
% Empirical Comparison &
% Human vs. LLM (GPT-4) Code &
% LLM-generated code exhibits higher cyclomatic complexity; may be less maintainable. \\
% \hline

% Molison et al. 2025 \cite{Molison2025} &
% Empirical Comparison &
% Human vs. LLM Code &
% LLM-code has fewer simple bugs but introduces more critical issues in complex scenarios. \\
% \hline

% Yang et al. 2024 \cite{Yang2024} &
% Empirical Study &
% LLM (Open-Source, GPT-4) vs. Evosuite &
% LLMs can outperform traditional tools, but performance is highly dependent on prompt design; limitations remain. \\
% \hline

% Beller et al. 2015 \cite{Beller2015} &
% Empirical Study &
% Test Assertions vs. Mutation Score &
% Assertion quantity and type are strongly correlated with test effectiveness (Nuance: assertTrue > assertEquals). \\
% \hline

% Peruma et al. \cite{Peruma2020} &
% ML Prediction &
% Static Metrics \& Test Smells &
% Static analysis (AST features, smells) can be used to effectively predict test flakiness. \\
% \hline

% Bellen et al. \cite{Bellen2013} &
% Static Analysis &
% Test Smells &
% Demonstrates using AST analysis to detect test smells (anti-patterns) that are root causes of flakiness. \\
% \hline

% \end{tabular}
% \end{table*}

\section{Study Design}
\label{Section:Methodology} 
This research conducts a large-scale comparative study of software tests written by human developers and tests generated by AI agents. The study uses the AIDev dataset as its main data source because it contains real-world software development activity \cite{li2025riseaiteammatessoftware}. However, the dataset includes only metadata and code patches, so we retrieve the full source code of the test files from their original repositories. Our method then collects, processes, and analyzes these files in a systematic way. We evaluate test quality in three aspects: assertion strength and specificity, edge-case coverage, and test flakiness. Figure~\ref{fig:approach} shows the full workflow, which is described in detail in the following steps.

\paragraph{\textbf{Phase 1: Metadata Manifest Generation}}
We begin with the AIDev dataset, using the \texttt{all\_pull\_request.parquet} file. From this file, we define two groups for comparison: pull requests created by AI agents (Agent-PRs) and a matched control group of pull requests created by human developers (Human-PRs).We note that Agent-PRs may include human review or edits before merging. This creates a possible human-in-the-loop limitation in the study. For each pull request, we extract the metadata needed for repository-level data collection, including the full repository name, such as owner/repository, the pull request number, and the final commit SHA. This phase produces a Retrieval Manifest, which lists the artifacts to be collected from their original repositories.

\paragraph{\textbf{Phase 2: Test Artifact Discovery and Retrieval}}
This phase focuses on finding and collecting the test files needed for analysis. For each repository listed in the Phase 1 manifest, we clone a local copy and identify files that follow common test file naming patterns, such as \texttt{*\_test.py} and \texttt{Test*.java}. We then compare these files with the files modified in each target pull request, which we obtain through the GitHub API. This helps us identify the specific test files relevant to each pull request.

Next, we use the commit SHA from the manifest to retrieve the full source code of each selected test file \cite{li2025riseaiteammatessoftware}. This ensures that our analysis is based on the complete version of each file in the correct commit. After retrieval, each test file is parsed into an Abstract Syntax Tree (AST) for further analysis.

The file type distribution in the dataset is shown in Table~\ref{tab:filetype_dist}. Since Python files are the most common, this study focuses only on Python test files. Table~\ref{tab:general_stats} reports the general dataset statistics. In that table, A-Sliced refers to the first 24,941 agent files, while A-Sliced-Rand refers to a random subset of 24,941 agent files.

\begin{table}[t]
\centering
\caption{File type distribution across Agent and Human datasets.}
\label{tab:filetype_dist}
\begin{tabular}{lrr}
\toprule
\textbf{Extension} & \textbf{Agent (\%)} & \textbf{Human (\%)} \\
\midrule
py   & 35.73\% & 21.48\% \\
ts   & 17.74\% & 21.48\% \\
pyc  & 9.05\%  & --      \\
go   & 5.17\%  & 13.53\% \\
js   & 4.90\%  & 2.87\%  \\
tsx  & 3.44\%  & 5.72\%  \\
cpp  & 2.59\%  & --      \\
cc   & 2.01\%  & --      \\
json & 1.77\%  & 3.25\%  \\
png  & --      & 7.25\%  \\
pb   & --      & 4.28\%  \\
\midrule
\textbf{Total Files} & 707{,}744 & 135{,}676 \\
\bottomrule
\end{tabular}
\end{table}


\paragraph{\textbf{Phase 3: Multi-Dimensional Quality Analysis}}
In the third phase, we analyze the parsed ASTs using three separate pipelines. Each pipeline is connected to one of our main research questions.

\begin{itemize}
    \item \textbf{Pipeline A (Assertion Strength Analysis for RQ1)} This analyzes the ASTs of the test files to find and classify the assertion statements. Assertions are classified using a predefined taxonomy, such as weak assertions and strong assertions. This pipeline measures how meaningful and specific the validation checks are within the test files.
\end{itemize}

\textbf{White-Box Flow chart:}

\begin{enumerate}
    \item \textbf{Input:} The source code of a Python test file is used as input.

    \item \textbf{Parsing:} The \texttt{ast.parse()} function converts the source code into an Abstract Syntax Tree (AST).

    \item \textbf{Traversal:} The \texttt{TestAnalyzer.visit()} method starts from the root \texttt{Module} node and visits the nodes in the AST.

    \item \textbf{Identification:} While visiting the AST, the analyzer looks for \texttt{ast.Call} nodes, which represent function calls. It calls those inside test functions, where the function name starts with \texttt{test\_}.

    \item \textbf{Classification:} The analyzer extracts the assertion call, such as \texttt{self.assertEquals}, and identifies the assertion name, such as \texttt{assertEquals}. Then it compares this name with the predefined \texttt{WEAK\_ASSERTIONS} and \texttt{STRONG\_ASSERTIONS} sets.

    \item \textbf{Advanced Predicate Analysis:} In the refined taxonomy, calls to \texttt{self.assertTrue} are checked more carefully using \texttt{\_is\_strong\_assertTrue(node)}.
    \begin{itemize}
        \item If the argument is a simple variable, such as \texttt{assertTrue(variable)}, the assertion is classified as \textbf{Weak}.
        \item If the argument contains a comparison, binary operation, or another function call, such as \texttt{assertTrue(x > 5)}, the assertion is classified as \textbf{Strong} because it checks more specific program behavior.
    \end{itemize}

    \item \textbf{Aggregation:} The analyzer updates the \texttt{weak\_count} and \texttt{strong\_count} values based on the classification.

    \item \textbf{Output:} The final output is the percentage of weak assertions and the percentage of strong assertions in the test file.
\end{enumerate}

This refined taxonomy is important because assertion strength depends on what the assertion checks, not only on the assertion name. For example, a basic taxonomy may classify \texttt{assertTrue(complex\_validation\_function(x))} as weak only because it uses \texttt{assertTrue}. However, this assertion may still contain meaningful validation logic. White-box refinement reduces this kind of misclassifications and makes the analysis more accurate.

\begin{table*}[h!]
\centering
\caption{Taxonomy of assertion strength.}
\begin{tabular}{|p{3.5cm}|p{2.2cm}|p{4cm}|p{3cm}|p{4cm}|}
\hline
\textbf{Assertion Type} &
\textbf{Simple Taxonomy} &
\textbf{Justification (Simple)} &
\textbf{Advanced Taxonomy} &
\textbf{Justification (Advanced)} \\
\hline

\texttt{assertNotNull(x)} &
Weak &
Checks for existence, not value. &
Weak &
Confirmed; minimal semantic validation. \\
\hline

\texttt{assertTrue(x)} &
Weak &
Checks for truthiness, not value. &
Context-Dependent &
Predicate must be inspected. \\
\hline

\texttt{assertTrue(variable)} &
Weak &
(See above) &
Weak &
Argument is \texttt{ast.Name}; low-value check. \\
\hline

\texttt{assertTrue(x > 5)} &
Weak &
(See above) &
Strong &
Argument is \texttt{ast.Compare}; encodes specific domain logic. \\
\hline

\texttt{assertTrue(is\_valid(x))} &
Weak &
(See above) &
Strong &
Argument is \texttt{ast.Call}; encodes specific semantic logic. \\
\hline

\texttt{assertEquals(5, x)} &
Strong &
Checks for a specific, non-trivial value. &
Strong &
Confirmed; provides precise state validation. \\
\hline

\texttt{assertRaises(MyError)} &
Strong &
Verifies a specific error-handling path. &
Strong &
Confirmed; tests non-happy paths. \\
\hline

\end{tabular}
\end{table*}
    \begin{itemize}
   \item  \textbf{Pipeline B (Edge Case Coverage Analysis for RQ2)} also traverses the ASTs to identify literal values (e.g., \texttt{0}, \texttt{-1}, \texttt{null}, empty strings) used as inputs in method calls. These are matched against a heuristic list of common edge cases to produce metrics on test coverage.\\
    The "Edge-Case Coverage" pipeline (Pipeline B) evaluates the test's thoroughness by identifying inputs that represent common boundary conditions.\\
    \end{itemize}
    \textbf{White-Box Flowchart:}
    \begin{enumerate}
    \item \textbf{Input:} The pipeline begins with the AST \texttt{Module} node.
    \item \textbf{Traversal:}  
    All function call nodes inside test methods are intercepted via  
    \texttt{TestAnalyzer.visit\_Call(node)}.
    \item \textbf{Argument Inspection:}  
    For each function call, the system iterates through every argument in  
    \texttt{node.args}.
    \item \textbf{Literal Matching:}
    \begin{itemize}
    \item Check whether \texttt{arg} is an \texttt{ast.Constant}.
     \item If true, compare \texttt{arg.value} against the heuristic list:
     \texttt{None}, \texttt{0}, \texttt{-1}, \texttt{""}.
    \end{itemize}

    \item \textbf{Collection Matching:}
    \begin{itemize}
    \item Check whether the argument is an \texttt{ast.List} or \texttt{ast.Dict}.
     \item If so, determine whether the corresponding element list (e.g., \texttt{arg.elts}) is empty.
    \end{itemize}
    The final metric \texttt{variety\_count} is computed as the size of this set(i.e., \texttt{len(categories)}), representing the number of distinct  
    edge-case types covered in the file.

    \item \textbf{Aggregation:}  
    A set is used to accumulate unique edge-case categories identified, such as:  
    \texttt{NULL\_INPUT}, \texttt{ZERO\_INPUT}, \texttt{EMPTY\_COLLECTION}.

    \item \textbf{Output:}  
    The final metric \texttt{variety\_count} is computed as the size of this set  
    (i.e., \texttt{len(categories)}), representing the number of distinct  
    edge-case types covered in the file.
\end{enumerate}

This static, literal-based approach is a heuristic. Its primary limitation, which must be This static, literal-based approach is a heuristic. Its primary limitation,is that it only detects edge cases involving literals (e.g., 0, None, ''), and cannot detect variable-driven edge cases (e.g., \texttt{empty\_list = [...]}; \texttt{myfunc(empty\_list))}, as this would require complex data-flow analysis, which is beyond the scope of this static AST analysis.

\begin{table*}[h!]
\centering
\caption{Heuristic Rules for Identifying Edge-Case Literals (RQ2).}
\begin{tabular}{|p{3.5cm}|p{3.5cm}|p{5cm}|p{5cm}|}
\hline
\textbf{Edge-Case Category} &
\textbf{Python Literal Value} &
\textbf{Target AST Node} &
\textbf{Example Code Fragment} \\
\hline

Null Inputs &
\texttt{None} &
\texttt{ast.Constant(value=None)} &
\texttt{my\_func(None)} \\
\hline

Empty Collections &
\texttt{[]} or \texttt{\{\}} &
\texttt{ast.List(elts=[])} or \texttt{ast.Dict(keys=[])} &
\texttt{my\_func([])} \\
\hline

Empty Strings &
\texttt{""} or \texttt{''} &
\texttt{ast.Constant(value="")} &
\texttt{my\_func("")} \\
\hline

Numeric Boundaries (Zero) &
\texttt{0} &
\texttt{ast.Constant(value=0)} &
\texttt{my\_func(0)} \\
\hline

Numeric Boundaries (Negative) &
\texttt{-1} &
\texttt{ast.Constant(value=-1)} &
\texttt{my\_func(-1)} \\
\hline

\end{tabular}
\end{table*}

\begin{itemize}
    \item \textbf{Pipeline C (Flakiness Analysis for RQ3)} uses a hybrid approach to identify possible flaky tests. First, it uses static analysis to detect common patterns that may lead to flaky behavior, such as \texttt{time.sleep}, network calls, file I/O, or non-deterministic functions such as \texttt{random()}. Then, for a sampled subset of tests that contain these patterns, we use dynamic analysis by running each test multiple times in an isolated environment. For example, each test may be executed 100 times to check whether the result changes across runs. The final output is a comparative flakiness rate for each group.
\end{itemize}

\textbf{White-Box Flowchart:}

\begin{enumerate}
    \item \textbf{Input:} The input is the AST \texttt{Module} node of a Python test file.

    \item \textbf{Traversal:} The \texttt{TestAnalyzer.visit\_Call(node)} method visits function calls inside test methods.

    \item \textbf{Path Extraction:} The helper function \texttt{\_get\_call\_path(node)} examines \texttt{node.func} and extracts the full call path, such as \texttt{('time', 'sleep')} or \texttt{('open',)}.

    \item \textbf{Pattern Matching:} The extracted call path is compared with the \texttt{FLAKINESS\_INDICATORS} dictionary. For example:
    \begin{itemize}
        \item \texttt{('time', 'sleep')} maps to \texttt{ASYNC\_WAIT}.
        \item \texttt{('random', 'random')} maps to \texttt{NON\_DETERMINISM}.
        \item \texttt{('open',)} maps to \texttt{FILE\_IO}.
    \end{itemize}

    \item \textbf{Aggregation:} Each matched indicator is added to the \texttt{rq3 \_flakiness \_indicators} list. The stored information also includes the line number, \texttt{node.lineno}, so that the matched code can be reviewed later.

    \item \textbf{Output:} The \textit{candidate\_rate} is calculated as \texttt{len(rq3\_flakiness\_indicators)}. This value represents the number of flakiness indicators found in the test file.
\end{enumerate}

This approach follows prior research on test flakiness and test smells. Such work often uses AST-based static analysis to detect code patterns, such as async waits or resource-related issues, that are common causes of flaky tests.

\begin{table*}[h!]
\centering
\caption{Rules for Identifying Flakiness Root Causes (RQ3).}
\begin{tabular}{|p{4cm}|p{4cm}|p{9.5cm}|}
\hline
\textbf{Flakiness Root Cause} &
\textbf{Code Pattern / Anti-Pattern} &
\textbf{Python AST Signature (Simplified)} \\
\hline

Async Wait &
Fixed time delays &
\texttt{ast.Call(func=ast.Attribute(value=ast.Name(id='time'), attr='sleep'))} \\
\hline

Non-Determinism &
Use of random functions &
\texttt{ast.Call(func=ast.Attribute(value=ast.Name(id='random'), attr='random'))} \\
\hline

Environment / Resource Leak &
Hard-coded File I/O &
\texttt{ast.Call(func=ast.Name(id='open'))} \\
\hline

Environment / Resource Leak &
Network Calls &
\texttt{ast.Call(func=ast.Attribute(value=ast.Name(id='socket'), attr='socket'))} \\
\hline

Non-Determinism &
System Time Dependency &
\texttt{ast.Call(func=ast.Attribute(value=ast.Name(id='datetime'), attr='now'))} \\
\hline

\end{tabular}
\end{table*}

\paragraph{\textbf{Phase 4: Synthesis and Comparative Evaluation}}
The final phase involves the synthesis of our findings. We aggregate metrics from all three analysis pipelines and perform rigorous statistical comparisons between the \textit{AgentTest-PRs} and \textit{HumanTest-PRs} cohorts for each quality dimension. These results are analyzed to directly answer our research questions, and the findings are synthesized into a final discussion, conclusions, and a set of actionable recommendations for both developers and AI tool builders.
%left bottom right top
% --- Figure 1: Approach Overview ---
\begin{figure*}[t]
    \centering
    \small % Use a smaller font for the diagram
    \includegraphics[width=15cm, height=15cm, trim=0 120 0 20, clip]{Images/Approach_Overview.png}
    \caption{Overview of our study design.}
    \label{fig:approach}
\end{figure*}

\subsection{``White-Boxing'' the Analysis Pipelines}
This section moves beyond the high-level overview to detail exactly how inputs are processed and features are calculated for each analysis pipeline. The implementation is a Python script utilizing the native AST (Abstract Syntax Trees) module.The script defines a single TestAnalyzer class that inherits from ast.NodeVisitor which traverse a Python test file's AST and collect metrics for all three research questions simultaneously.

\subsection{Manual Analysis and Exploratory Investigation}
\label{ssec:explore and investigation}
Prior to embarking on a full-scale, automated analysis, we  undertake a foundational phase of manual, qualitative analysis. This will involve the detailed inspection of a small, randomly selected sample of 50 pull requests from each of our two cohorts (Agent-authored and Human-authored). This exploratory stage is deliberately structured to follow the established ``Explore, Analyze, Synthesize'' framework for qualitative data analysis, allowing us to develop a nuanced, human-centric understanding of the data before applying automated techniques. %During this phase, we will manually inspect the test files with the following objectives:

% \paragraph{\textbf{Explore}}
% To cultivate a qualitative and intuitive understanding of the typical structure, stylistic conventions, and overall complexity that characterize tests written by AI agents in comparison to those crafted by their human counterparts.
% \paragraph{\textbf{Analyze}}
% Identify and catalog recurring patterns, such as the frequent use of specific assertion types,the presence or absence of explanatory comments detailing test logic, or common approaches to test setup and tear down.The insights gained from this pattern analysis will be instrumental in refining the taxonomies and classification schemes to be used in our subsequent, large-scale automated analysis.
% \paragraph{\textbf{Synthesize}}
% Formulate a set of initial data-based hypotheses regarding the potential differences in test quality between the two cohorts. These preliminary hypotheses will serve as a guiding framework for the interpretation of our quantitative results and will help to contextualize the statistical findings within a richer and more descriptive narrative.

\subsection{Automated Analysis Pipelines}
\label{ssec:data-analysis}
The analytical core of our research is made up of three different, automated analysis pipelines, which integrate both static and dynamic techniques to systematically quantify the multifaceted dimensions of test quality.

\paragraph{\textbf{Assertion Strength Analysis}}
To measure the semantic depth of the tests, we use tree-sitter to generate an Abstract Syntax Tree (AST) for each collected test file. We then traverse each AST to identify method calls that correspond to assertion functions in common testing frameworks, such as \texttt{assertEquals}, \texttt{assertTrue}, and \texttt{assertIn}.

Each assertion is classified using a predefined strength taxonomy. Assertions that check only basic conditions, such as non-null values using \texttt{assertNotNull} or simple truth values using \texttt{assertTrue(variable)}, are classified as \textbf{Weak}. These assertions provide limited information about whether the program behavior is correct.

In contrast, assertions that check specific expected values, exact exception types, or deep equality between complex data structures are classified as \textbf{Strong}. Examples include \texttt{assertEquals(expected, actual)} and assertions that compare nested objects or collections. These assertions provide a stronger form of validation because they check more precise program behavior.
    
\paragraph{\textbf{Edge-Case Coverage Analysis}}
Using the same set of ASTs, this pipeline evaluates how well the tests cover edge cases and boundary inputs. The pipeline identifies method calls inside test functions and extracts the literal values passed as arguments. These values may include numbers, strings, and named constants.

The extracted values are then compared with a predefined list of common edge cases and boundary values. This list includes values such as \texttt{0}, \texttt{-1}, \texttt{None}, and \texttt{""}, which represents an empty string. It may also include language-specific boundary values such as \texttt{Integer.MAX\_VALUE} and \texttt{Integer.MIN\_VALUE}.

We use the frequency and variety of these edge-case values as a quantitative measure of test robustness. A test file with more diverse edge-case inputs is considered more likely to check behavior near the boundaries of the input domain.

\paragraph{\textbf{Flakiness Analysis}}
This analysis uses a two-stage process. In the first stage, we perform static analysis on the test source code to identify patterns that are commonly linked to flaky tests. These patterns include the use of non-deterministic functions, such as \texttt{random()}, hard-coded time delays, such as \texttt{Thread.sleep}, interactions with networks or file systems, and dependencies on system time.

Any test that contains one or more of these patterns is marked as a \textit{flakiness candidate}. In the second stage, we perform dynamic analysis on a representative sample of these candidates. For each selected test, we create an isolated execution environment, check the specific software commit linked to the test, and run the test repeatedly, such as \texttt{N = 100} times.

A test is confirmed to be flaky if it produces both passing and failing results across these repeated runs without any changes to the source code.

 
 \section{Experimental Results}
\label{Section:Result}
To rigorously evaluate the qualitative and quantitative characteristics of agent-generated software tests, our experimental framework is structured around three central research questions (RQs). Each of these questions is designed to probe one of the critical dimensions of test quality that were identified in our study design. %The empirical basis for our experiments will be a direct and systematic comparison of the analytical results obtained \hl{from the}

The following section %replace the "Expected Results" from the Phase 2 study design with the actual results 
 derived from executing the TestAnalyzer script on the two
\textit{AgentTest-PRs} cohort against the established baseline provided by the cohorts of 20 test files each. The raw, aggregated data was sourced from the script's JSON outputs.
\textit{HumanTest-PRs} cohort.

\subsection{RQ1: How does the assertion strength of agent-generated tests compare to human-written tests?}
\label{ssec:rq1}
 
\paragraph{\textbf{Methodology}}
For each test file in both cohorts, we classified every assertion as either ``Weak'' or ``Strong'' using the taxonomy in Table~\ref{tab:assertion-taxonomy-actual}. We then computed the weak-to-total ratio per file. Files where this ratio is high are considered to have shallow assertion depth. Aggregating these per-file ratios gives us a distribution for each cohort, which we compare directly.

\paragraph{\textbf{Weak Assertions}}
We treat an assertion as weak if it only checks whether something exists, whether a boolean condition holds, or whether an object is of a certain type, without pinning down the actual value or internal state. Typical examples are \texttt{assertNotNull}, \texttt{assertTrue}, and \texttt{assertIsInstance}~\cite{schafer2023empirical}.

\paragraph{\textbf{Strong Assertions}}
A strong assertion nails down a specific expected output, confirms a particular state change, or checks a non-trivial behavioral property. Think \texttt{assertEquals(5, result)}, \texttt{assertThrows(SpecificException.class)}, or \texttt{assertDeepEquals(expectedObject, resultObject)}~\cite{schafer2023empirical}.

\begin{table}[t]
\centering
\caption{Assertion categories used in our analysis.}
\begin{tabular}{p{2.1cm} p{5.9cm}}
\toprule
\textbf{Category} & \textbf{Assertion Methods} \\
\midrule

Weak &
\texttt{assertTrue},
\texttt{assertFalse},
\texttt{assertIs},
\texttt{assertIsNot},
\texttt{assertIsNone},
\texttt{assertIsNotNone},
\texttt{assertIn},
\texttt{assertNotIn},
\texttt{assertIsInstance},
\texttt{assertNotIsInstance} \\

\midrule

Strong &
\texttt{assertEquals},
\texttt{assertNotEquals},
\texttt{assertEqual},
\texttt{assertNotEqual},
\texttt{assertDictEqual},
\texttt{assertListEqual},
\texttt{assertSetEqual},
\texttt{assertTupleEqual},
\texttt{assertRaises},
\texttt{assertRaisesRegex},
\texttt{assertLogs},
\texttt{assertGreater},
\texttt{assertGreaterEqual},
\texttt{assertLess},
\texttt{assertLessEqual},
\texttt{assertAlmostEqual},
\texttt{assertNotAlmostEqual} \\

\midrule

Reclassified (Advanced) &
\texttt{assertTrue(expr)} with logical/comparison or function call expressions (e.g., \texttt{is\_valid(x)}, \texttt{x > 0}). \\

\midrule

Unknown &
Project-specific or misspelled assertions, e.g., \texttt{assertSomething}, \texttt{assertIsLess}, or bare \texttt{assert}. \\

\bottomrule
\end{tabular}
\label{tab:assertion-taxonomy-actual}
\end{table}




\paragraph{\textbf{Metrics}}
\begin{itemize}
    \item \textbf{Percentage of Weak Assertions:}
    (Count of Weak Assertions / Total Count of Assertions) $\times$ 100
    \item \textbf{Percentage of Strong Assertions:}
    (Count of Strong Assertions / Total Count of Assertions) $\times$ 100
\end{itemize}

\paragraph{\textbf{Results}}
The numbers in Table~\ref{tab:general_stats} turned out closer than we expected. Going in, we assumed there would be a clear quality gap between human and agent assertions, but that is not quite what happened. Agents did lean more toward weak assertions (14.63\% vs.\ 11.92\% for humans), and humans held a corresponding edge in strong ones (88.08\% vs.\ 85.37\%). But this is not the large divide that earlier work on semantic comprehension issues in LLM-generated code might lead you to expect~\cite{schafer2023empirical}.

The more telling result is in the ``Unknown'' column. Agents produced unrecognized assertion patterns at 11.58\%, roughly eight times the 1.46\% we saw in human-written tests. What this means in day-to-day terms is that agents frequently reach for assertion methods that either do not exist in standard libraries or are misspelled. A reviewer looking at one of these test files has to pause and work out whether the assertion is doing something clever with a project-specific helper, or whether the agent just made it up. Across a large codebase, that kind of friction adds up and can quietly erode maintainability.
 
The takeaway here is that agents are not dramatically worse than humans in raw assertion strength. But the gap in discipline, writing assertions that a reader can immediately trust, matters more than the percentages alone would suggest.

\begin{table}[t]
\centering
\scriptsize
\caption{General statistics for human vs.\ agent cohorts.}
\label{tab:general_stats}
\begin{tabular}{lrrrr}
\toprule
Metric & Human & Agent & A-Sliced & A-Sliced-Rand \\
\midrule
Total Files & 24{,}941 & 179{,}732 & 24{,}941 & 24{,}941 \\
Files w/Asserts & 1{,}730 & 29{,}353 & 3{,}412 & 4{,}046 \\
Weak (\%) & 11.92 & 14.30 & 13.90 & 14.63 \\
Strong (\%) & 88.08 & 85.70 & 86.10 & 85.37 \\
Unknown (\%) & 1.46 & 10.93 & 9.22 & 11.58 \\
Cand.\ Rate & 0.30 & 0.44 & 0.46 & 0.41 \\
Variety & 0.32 & 0.61 & 0.58 & 0.62 \\
\bottomrule
\end{tabular}
\end{table}


\begin{figure}[h]
    \centering
    \includegraphics[width=\linewidth]{Images/RQ1_new.jpg}
    \caption{Assertion strength comparison across cohorts. Agents exhibit a higher rate of `Unknown' assertion types than humans.}
    \label{fig:rq1_assertion_strength}
\end{figure}

\subsection{RQ2: To what extent do agent-generated tests cover edge and boundary cases compared to human-written tests?}
\label{ssec:rq2}
 
\paragraph{\textbf{Methodology}}
We scanned the input values used in method calls across all collected test files. Our pipeline flags literal values and constants that match known boundary conditions and special cases: zero, negative numbers, null values, empty collections, and numeric or string length limits. For each test file, we recorded how many of these categories appeared.
 
\paragraph{\textbf{Metrics}}
\begin{itemize}
    \item {\textbf{Edge Case Frequencies.}}
    The average number of distinct edge case patterns found per 1,000 lines of test code. This gives us a normalized density that is comparable across files of different sizes.
    \item {\textbf{Edge Case Variety.}}
    The percentage of tests in each cohort that include at least one test targeting our predefined categories: null inputs, empty collections, numeric boundaries (zero, negative values, maximum integers), and strings of zero or very long length.
\end{itemize}
 

\paragraph{\textbf{Results}}
We expected humans to do better here, but the data in Table~\ref{tab:edge_counts} says otherwise. Agents beat humans on nearly every edge case category. For ``Null Input,'' agents hit 13.4\% compared to 8.3\% for humans. ``Empty Collection'' showed a similar gap: 14.9\% vs.\ 8.1\%. The biggest difference was in ``Zero Input,'' where agents reached 27.7\%, more than double the human rate of 11.2\%.
 
The variety scores tell a similar story. Agents averaged 0.614 on our variety metric, almost twice the human score of 0.318 (Table~\ref{tab:general_stats}). So when agents write tests, they tend to throw a wider net of input types into a single suite. One thing both groups got wrong: neither humans nor agents tested negative numeric boundaries at all (0.0\% across the board), which is a blind spot worth noting.
 
These results push back against the idea that LLMs are bad at reasoning about boundary conditions. What seems to be happening instead is that agents default to listing out standard boundary values (null, zero, empty) almost mechanically, while human developers are more likely to skip cases they consider unlikely based on their understanding of the code. Whether the agent's extra coverage is genuinely useful or just noise is a separate question, but the raw numbers clearly favor agents on this metric.

\begin{table}[t]
\centering
\scriptsize
\caption{Prevalence of edge cases in assertions.}
\label{tab:edge_counts}
\begin{tabular}{lrrrr}
\toprule
Edge Case & Human & Agent & A-Sliced & A-Sliced-Rand \\
\midrule
Null Input        & 8.3\%  & 13.0\% & 13.5\% & 13.4\%  \\
Zero Input        & 11.2\% & 27.7\% & 23.7\% & 27.7\%  \\
Negative Input    & 0.0\%  & 0.0\%  & 0.0\%  & 0.0\%   \\
Empty String      & 4.1\%  & 5.7\%  & 5.4\%  & 5.8\%   \\
Empty Collection  & 8.1\%  & 15.0\% & 15.3\% & 14.9\%  \\
\bottomrule
\end{tabular}
\end{table}


\begin{figure}[h]
    \centering
    \includegraphics[width=\linewidth]{Images/RQ2_new.jpg}
    \caption{Edge-Case coverage comparison. Agents significantly outperform humans in covering Zero Inputs and Null Inputs.}
    \label{fig:rq2_edge_coverage}
\end{figure}

\subsection{RQ3: What is the prevalence of flaky tests in pull requests submitted by AI agents compared to those submitted by humans?}
\paragraph{\textbf{Methodology}}
Our two-stage flakiness analysis will first use static analysis to find a complete list of ``flakiness candidates'' from both the agent- and human-authored groups.  This identification relies on the existence of recognized anti-patterns and code constructs that exhibit a strong correlation with non-deterministic test behavior, including the utilization of Thread.sleep, interactions with network endpoints, file I/O operations, and invocations of non-deterministic APIs.  After this static screening, we move on to a dynamic analysis phase. In this phase, we will run a random sample of 1,000 tests from each group of candidates.  We will run each of these tests 100 times in a controlled, separate space.  If a test shows both a passing and a failing result at least once during these repeated runs, it will be clear that it is flaky.

\paragraph{\textbf{Metrics}}
\begin{itemize}
    \item \textbf{Candidate Rate:} The percentage of tests in a cohort that our static analysis flags as potential flakiness candidates.
    \item \textbf{Confirmed Flakiness Rate:} Of the sampled candidates, the percentage that actually turned out to be flaky when run through our dynamic analysis.
\end{itemize}
 
\paragraph{\textbf{Results}}
This is the one research question where our initial guess held up. Agents do appear more prone to introducing instability. Their average candidate rate came in at 0.435, well above the human rate of 0.301 (Table~\ref{tab:general_stats}). Looking at what drives that gap in Table~\ref{tab:type_counts}, agents leaned harder on ``Non-Determinism'' (5.2\% vs.\ 3.1\%) and ``File I/O'' (4.4\% vs.\ 3.5\%). Both of these point to agents not being careful about side effects and resource isolation when generating test code.
 
The concurrency picture is less clear-cut. Humans actually used ``Async Wait'' patterns slightly more (1.9\%) than agents did (1.4\%), and network I/O was negligible for both groups. So it is not that agents are worse across every flakiness dimension. The problem is narrower than that: agents are specifically bad at handling random values and file system access. They write code that touches the disk or calls non-deterministic APIs without mocking or cleanup, and that is where most of the extra flakiness risk comes from. Humans write flaky code too, of course, but agents do it at a noticeably higher rate in these two specific areas.

\begin{table}[t]
\centering
\scriptsize
\caption{Prevalence of specific code patterns.}
\label{tab:type_counts}
\begin{tabular}{lrrrr}
\toprule
Category & Human & Agent & A-Sliced & A-Sliced-Rand \\
\midrule
Async Wait       & 1.9\% & 1.3\% & 1.4\% & 1.4\% \\
Non-Determinism  & 3.1\% & 5.2\% & 5.8\% & 5.2\% \\
File I/O         & 3.5\% & 4.6\% & 3.9\% & 4.4\% \\
Network I/O      & 0.1\% & 0.2\% & 0.2\% & 0.1\% \\
\bottomrule
\end{tabular}
\end{table}


\begin{figure}[h]
    \centering
    \includegraphics[width=\linewidth]{Images/RQ3_new.jpg}
    \caption{Flakiness indicator rates. Agents exhibit higher rates of Non-Determinism and File I/O usage.}
    \label{fig:rq3_flakiness}
\end{figure}

% \section{Discussion}
% \label{Section:Discussion}


% The results of this study present a complex picture of the current state of agent-based test generation. While our initial hypotheses anticipated that human developers would outperform AI agents across all dimensions of test quality, the empirical data reveals a distinct trade-off between \textit{coverage breadth} and \textit{execution stability}.

% \subsection{Concrete Example: Human vs.\ Agent PR}
% As a small Phase-4 extension, we briefly describe one anonymized pair of PRs we inspected manually:

% \begin{itemize}
%   \item \textbf{Human PR:} Added a new validation function for user profiles and wrote three tests: a normal case, a missing required field case, and a malformed email case. All assertions were strong and the tests used fixed example inputs.
%   \item \textbf{Agent PR:} In a different repository, an AI agent added a similar validation helper. It wrote tests for normal, missing field, null profile, empty name, very long name, and a handful of random strings. Some tests used only weak assertions (checking that ``no exception is thrown''), and one test created a temp file to log output, which we flagged as a flakiness indicator.
% \end{itemize}
% \begin{figure}[t]
% \centering
% \begin{subfigure}{0.48\columnwidth}
%   \centering
%   \includegraphics[width=\linewidth]{Images/human_test_continuous_profiler.png}
%   \caption{Human test file (sentry-python).}
%   \label{fig:human-test-example}
% \end{subfigure}
% \hfill
% \begin{subfigure}{0.48\columnwidth}
%   \centering
%   \includegraphics[width=\linewidth]{Images/agent_test_alert_notifications.png}
%   \caption{Agent-generated test file (Gal-Friday2).}
%   \label{fig:agent-test-example}
% \end{subfigure}
% \caption{Concrete example of a human-written test (left) and an agent-generated test (right). The human file focuses on a specific profiling feature, whereas the agent file exercises alert notifications with a broader but noisier set of behaviors.}
% \label{fig:human-vs-agent-code}
% \end{figure}

% The qualitative comparison in Figure~\ref{fig:human-vs-agent-code} shows the same trend observed in our quantitative metrics: human tests are concise and outcome-oriented, whereas agent tests tend to be broad, exploratory, and occasionally noisy.

% The human test suite was small and focused; the agent suite was broad but noisy. In an ideal workflow, the human team would merge a cleaned-up subset of the agent tests, preserving the useful edge cases while trimming the risky patterns.


% \subsection{The Edge Case Paradox: Exhaustive Enumeration vs. Contextual Pragmatism}
% Perhaps the most striking finding of this study is the inversion of our expectations regarding edge case coverage (RQ2). Contrary to the belief that agents lack the reasoning to identify boundary conditions, AI agents significantly outperformed humans in generating tests for inputs such as zero, null values, and empty collections. 

% We hypothesize that this stems from the fundamental nature of Large Language Models as pattern-matching engines. During training, these models ingest vast quantities of unit testing boilerplate where "null checks" and "zero checks" are ubiquitous patterns. Consequently, agents appear to default to an \textit{exhaustive enumeration} strategy, mechanically generating tests for standard boundary values regardless of the specific business context. In contrast, human developers appear to exercise \textit{contextual pragmatism}; a human developer might skip a "zero input" test if they know the upstream logic makes such a state impossible. While the agent's approach yields higher variety scores (0.614 vs. 0.318), it raises questions about the utility of these tests—whether they represent genuine robustness or merely "testing theater."
% \subsection{Assertion Quality and the "Unknown" Drift}
% In terms of assertion strength (RQ1), the gap between humans and agents was narrower than expected (88.08\% vs 85.37\% strong assertions). However, the high prevalence of "Unknown" assertion types in agent-generated code (11.58\%) warrants close attention. This statistic suggests that while agents can mimic the structure of strong tests, they frequently veer into non-standard validation logic or hallucinated assertion libraries. This "drift" necessitates a higher cognitive load for human reviewers, who must decipher whether an unfamiliar assertion pattern is a clever validation technique or a subtle hallucination.
% \subsection{Reliability and Environmental Awareness}
% The flakiness analysis (RQ3) highlights the most significant current limitation of agent-generated tests: a lack of environmental awareness. Agents demonstrated a significantly higher candidate rate for flakiness, driven largely by Non-Determinism and File I/O. 

% This suggests that while agents understand the \textit{syntax} of testing, they struggle with the \textit{semantics} of test isolation. A human developer implicitly understands that reading from a real file system introduces fragility and will likely implement mocking strategies. An agent, lacking a mental model of the CI/CD environment, is more likely to generate code that performs direct I/O or relies on unstable random seeds. This finding implies that currently, agent-generated tests require a sandbox or rigorous linting pipeline (including static and dynamic analysis) to detect and remove non-determinisitic patterns (such as file IO or randomization) before integration. Humans display an implicit advantage in isolating environmental dependencies and maintaining test stability.

% \subsection{Implications for Software Engineering}
% The data suggests that AI agents are shifting from being "code completion" tools to "fuzzing" partners. They excel at identifying generic failure modes (nulls, zeros) that humans frequently overlook due to fatigue or overconfidence. However, their propensity for introducing flakiness and non-standard assertions indicates that they are not yet ready for autonomous test suite management. The optimal workflow is a hybrid model: using AI agents to generate broad edge-case and test scaffolds, then requiring human review to refine assertions and mock desperation for reliability, This approach leverages the breadth of AI while maintaining human judgment for depth and stability.

\section{Threats To Validity}
\label{Section:Threats}
%A rigorous component of any empirical study is the acknowledgment of it's limitations. This section discusses threats to validity using the standard four-category framework: conclusion, internal, construct, and external validity.\cite{chowdhury2024swebenchverified}.These threats are not admissions of flawed design, but rather a necessary contextualization of the results and a guide for future research.
\subsection{Threats to Conclusion Validity}
The validity of the conclusion concerns the ability to draw a correct conclusion about the treatment (author type) and the outcome (test quality metrics).
    \begin{itemize}
        \item \textbf{Statistical Power and Sample Size:} The study's comparative analysis depends on statistical significance. The ability to draw a statistically significant conclusion (e.g. ``AI test has a higher proportion of weak assertions'') depends on the final sample size (number of test files) and the variance within each cohort. The sample size used in the experimental results (N=29491 per cohort) is not enough, and the variance in some cases is low (e.g., RQ3). A larger sample is required to draw more definitive conclusions.
        \item \textbf{Reliability of Measures:} The analysis relies on the automated AST pipelines. A bug in the TestAnalyzer (e.g., misclassifying an assertion) would systematically skew the results and threaten then conclusion. This threat is mitigated by the ``white-box'' design, which makes all heuristics explicit and renewable.
    \end{itemize}

% \subsection{Threats to Internal Validity:}
% Internal validity concerns confounding variables that could provide an alternative explanation for the observed results.
%     \begin{itemize}
%         \item \textbf{The "Junior Collaborator" Confounding:} This is the most significant threat to the study's internal validity. The research is framed as a comparison between "AI" and "Human" authors. However, "Agent - PRs" involve human reviewers who may edit or approve AI-generated tests before merging. This introduces a critical confounding variable: human developer experience.An observed "weak" test in the Agent-PR cohort could be the result of (a) the AI generating a weak test, or (b) a junior developer accepting a weak test from the AI, whereas a senior developer would have rejected it.Therefore, the study may not be measuring "AI vs. Human" but rather "Human-who-uses-AI" vs. "Human-who-does-not-use-AI".
%         \item \textbf{Data Sourcing Fidelity:} The Phase 2 artifact retrieval process is complex. An error in matching a pull request from the AIDev dataset to the correct commit SHA, or an error in the filters used to identify test files, could introduce noise or bias into the two cohorts, threatening the causal link between author type and test quality.
%     \end{itemize}
\subsection{Threats to Construct Validity:}
Construct validity concerns whether the metrics being measured (RQs 1-3) are valid proxies for the construct they are intended to represent (``test quality'').
\begin{itemize}
    \item \textbf{RQ1 (Assertion Strength):} The ``Advanced Taxonomy'' is a significant mitigation for the weak \texttt{assertTrue} construct. However, even this advanced proxy is limited. The true strength of assertTrue(is\_valid(x)) depends on the semantic complexity of the is\_valid(x) function, which the AST parser does not analyze. The metric is a proxy for strength, not a direct measure of it. This threat, particularly for RQ1, was confirmed to be the central finding of the experimental analysis. As discussed in earlier section, the construct (the TestAnalyzer) was a poor proxy for ``test quality'' in the Human-PR cohort due to its inability to parse pytest-style assertions, thereby invalidating a direct comparison.
    \item \textbf{RQ2 (Edge-Case Coverage):} The pipeline used in this study only finds literal constants (e.g., \texttt{myfunc(0)}). It cannot detect variables-based or context-dependent edge cases, since data-flow analysis is not performed. Therefore, results may underestimate edge case coverage, particularly in more advanced or human-authored tests. A test file could have 100\% comprehensive edge-case coverage using variables, and our tool would incorrectly score it at 0. Therefore, this metric is a proxy for ``developer awareness of literal edge cases,'' not ``true edge case coverage''.
    \item \textbf{RQ3 (Test Flakiness):} The construct validity for RQ3 is stronger, but only because the methodology is careful to define it. The static analysis (Pipeline C) does \textit{not} measure ``flakiness''; it measures ``flakiness \textit{candidates}''. A call to Thread.sleep, for example, is a known anti-pattern but can be a \textit{legitimate} part of a test (e.g., testing a timeout). The static metric \textit{alone} is a weak construct. The construct is only valid because it is the first stage of a two-stage (static + dynamic) process defined in the methodology.
\end{itemize}
\subsection{Threats to External Validity:}
External validity concerns the ability to generalize the study's findings to other contexts.
\begin{itemize}
    \item \textbf{Language-Specific Threat:} As detailed in earlier Sections, the provided implementation uses Python's AST module, Therefore, the findings will \textit{only} be valid for Python-based projects. These results may not carry over to Java, JavaScript, or C\# projects, since those languages come with their own testing frameworks, coding conventions, and AI generation behaviors. The RQ1 findings reinforced this concern. The choice of Python, with its idiomatic split between unit test and Pytest, was a decisive factor. A study in a language with a more unified testing framework (e.g., Java/JUnit) might yield different results.
    \item \textbf{Project-Specific Threat:} The AIDev dataset and the study's recovery methodology are based on open-source projects. The developers, motivations, the quality assurance processes in open-source software may not be representative of closed-source, industrial, or enterprise environments. The findings may not be generalized beyond the open-source context.
    \item \textbf{AI-Specific Threat:} The study analyzes ``Agent-PRs'' from the AIDev dataset. The quality of these PRs depends on the \textit{specific} AI models prevalent during the dataset's collection (e.g., older versions of GPT) and on the prompts provided by their human collaborators. The findings for ``AI'' may not be generalizable to future, more advanced models or to different prompting-strategy environments.
\end{itemize}
\section{Conclusion}
\label{Section:Conclusion}
To assess assertion strength, edge-case coverage, and flakiness, we analyzed more than 200,000 files in a comprehensive empirical comparison of human-authored and agent-generated test suites. What we found is less a story of one side being better, and more a trade-off: agents cover more ground mechanically, but they are less stable when it comes to execution environments. The common assumption that AI agents simply write worse code than humans does not hold up cleanly.
For assertion quality (RQ1), agents roughly matched humans in the split between strong and weak assertions. Where they fell short was in what we call ``assertion drift'': a high rate of unfamiliar or non-standard validation patterns that do not belong to any recognized testing library. This implies that agents can lack the semantic accuracy necessary for reliable verification, even when they can replicate the structure of rigorous testing.

Surprisingly, our examination of edge situations (RQ2) showed that AI greatly exceeds humans in terms of coverage frequency, especially for common boundary conditions such as zero and null inputs. This suggests that LLMs' probabilistic structure serves as a useful ``fuzzing'' technique by thoroughly listing cases that human developers frequently ignore because of exhaustion or implicit domain assumptions.

However, the results of our flaking investigation offset this benefit (RQ3). Non-deterministic behaviors were statistically more likely to be introduced by agent-generated tests, particularly through incorrect file input/output and uncontrolled random seed usage. The greatest obstacle to independent implementation of agent-based testing is still this lack of environmental knowledge.

In the end, we determine that AI agents now work best as high-volume ``test generators'' that need human supervision rather than as human testers' substitutes. They struggle with the depth of reliability (stability), they excel at broadening the scope of the test suite (coverage). Future research should focus on hybrid workflows in which human developers—or specialized static analysis tools—enforce semantic correctness and environmental isolation while agents suggest edge cases and assertion skeletons.

\section{Future Work}
\label{Section:FutureWork}
The results of this study highlight a number of crucial avenues for the development of automated tests in the future. Although the coverage breadths of current AI agents are impressive, the following areas require focused research due to their lack of environmental awareness and assertion precision.

\subsection{Environmental Modeling and Sandboxed Execution}
According to our data, the careless use of File I/O and non-deterministic APIs is a major cause of agent-generated flakiness (RQ3). Developing ``environment-aware'' agents that are specifically trained or prompted to identify the limitations of a CI/CD pipeline should be the main goal of future research. Integrating Sandboxed Execution Environments, where agents can run their generated tests iteratively, identify side effects (such as leftover files or thread leaks.) and self-correct before submission, is a promising approach. As a result, the paradigm changes from ``Generate and Pray'' to ``Generate, Execute, and Refine.''

\subsection{Hybrid Assertion Generation (RAG + Rule-Based)}
In order to address the high frequency of ``Unknown'' and weak assertions (RQ1), future research should investigate hybrid architectures that integrate static analysis rules with large language models. The target repository's specific assertion libraries and coding conventions could be dynamically fed to agents through the use of \textit{Retrieval-Augmented Generation (RAG)}. By grounding the agent's output in legitimate, project-specific syntax, this would lessen the ``assertion drift'' that our study found. Additionally, the semantic depth of generated tests could be greatly enhanced by fine-tuning models specifically on ``strong'' assertion patterns—those that verify state changes rather than just existence.

\subsection{Automated Mocking and Dependency Inference}
Agents are great at counting inputs, but have trouble with the context needed to test them safely, according to the ``Edge Case Paradox'' (RQ2). One area that deserves more attention is automated dependency inference. The idea is that agents would scan a method's call graph, figure out which external dependencies are involved (databases, APIs, time providers), and generate the appropriate mocks on their own. Agents would be able to take advantage of their edge-case creativity without adding the instability that comes with integration testing if they went beyond simple input generation to complete ``Test Fixture Synthesis.''

\subsection{Longitudinal Maintenance Studies}
Lastly, even though this study examined test development, it is still unknown how much maintenance of agent-generated suites will cost in the long run. Future long-term research should monitor the ``survival rate'' of agent-authored tests in industrial code bases, i.e., how frequently noise causes them to be removed, rewritten, or disabled. It is crucial to comprehend the ``Maintenance-to-Value'' ratio in order to ascertain the actual return on investment for software testing with AI.

\section{Data Availability} 
Our replication package is available online \cite{compsac_replication_package}.

\section{Acknowledgments} 
 During the preparation of this work, the authors used the ChatGPT Web interface to improve the language and readability, and used Gemini to create an approach overview figure . After using this tool, the authors reviewed and edited the content as needed and take full responsibility for the content of the publication. 

%\bibliographystyle{ACM-Reference-Format}
\bibliographystyle{abbrv}
\bibliography{sample-base}


\end{document}
\endinput
%%
%% End of file `sample-authordraft.tex'.
